\section{Discussion}
\label{sec:discussion}

\paragraph{The extraction contract is part of the measured system.}
The gap between strict and D-F1 exposes a practical distinction: a detector
can emit useful PII spans while failing the interface through which an
application consumes them. Strict rejection measures this deployment-facing
failure, whereas itemwise diagnostics expose the valid evidence remaining
in the same received response. The gap is not an additive estimate of
``performance lost to formatting'': the diagnostic changes both retained
predictions and false-positive accounting, and F1 is nonlinear. Reporting
both scores with compliance, empty-output rates and invalid-item counts
is more informative than selecting whichever score favors a system.

\paragraph{More context is an empirical choice.}
The ten complete pairs consistently favor local context on D-F1, but this
does not establish that financial attributes never require distant evidence.
Full and local conditions share the same explicit target, which may already
contain much of the information needed for the output task. Additional
context can be relevant without improving exact enumeration under a fixed
response budget. Attention dilution, instruction following and output
allocation are possible explanations, not mechanisms identified by these
experiments. Pointwise paired intervals quantify resampling variation over
the observed documents; they do not measure stochastic generation variance
or remove uncertainty about the serving environment.

\paragraph{Coverage matters alongside boundary accuracy.}
Strong identifier performance relative to attribute performance is consistent
with a limited operational coverage of the benchmark policy. It does not
show that contextual attributes are legally less important. Conversely,
an exact-boundary error on a long consultation clause can hide partial
recognition that would be visible under character overlap. The two views
answer different questions: whether the requested typed span was recovered,
and how much labeled text was covered. Character scores cannot charge a
character-level penalty to an unanchorable item and should therefore be
read with invalid-item counts and penalized exact precision. We do not
interpret overlap as a validated measure of remaining disclosure risk.

\paragraph{Implications for financial evaluation.}
These results support evaluating a detector as a policy, an input protocol
and an output contract together. A high format-compliance rate alone does
not show that sensitive information is protected; a low strict score alone
does not locate the failure. The benchmark makes these distinctions
inspectable using the same archived responses. For future deployments,
constrained generation or a different extraction interface would constitute
a new system configuration requiring its own prospective evaluation.
We do not retrofit such changes after observing this test set. The present
study supplies a reproducible characterization of the completed runs,
not a certification that any evaluated system is safe for customer records.
